\paragraph{C.3ay (Proved/refuted): critical restriction kernels in positive characteristic.}
For every prime $p$ and every $r\ge2$, the explicit construction below
gives an excellent normal strictly henselian local ring $R$ of dimension
$N=2r$, with algebraically closed coefficient and residue field of
characteristic $p$, and a nonzero divisor $f\in R$ such that both
punctures $X$ and $Y=X\cap V(f)$ are regular, while
\[
 \ker\!\left(\mathrm{CH}^r(X)_{\mathrm{alg}}
       \longrightarrow\mathrm{CH}^r(Y)_{\mathrm{alg}}\right)
                 \otimes\mathbb Q\ne0.
\]
The exhibited ordinary class has infinite order and its restriction is
zero integrally, also when $p$ divides $r$. The coefficient field is
the specified transcendental extension
$\overline{\overline{\mathbb F}_p(E_0)}$.

\begin{proof}
\emph{Step 1: the fixed elliptic curve and a specified non-torsion class.}
In plane coordinates $(U:V:W)$ take the elliptic curve $E_0$ over
$\mathbb F_p$ with origin $o=(0:1:0)$ and equation
\[
 \begin{cases}
 V^2W+VW^2+U^3=0,&p=2,\\
 V^2W-U^3+UW^2=0,&p\ne2.
 \end{cases}
\]
These are the all-prime cubics already checked in C.3x. For completeness,
in characteristic two their partial derivatives are $U^2,W^2,V^2$.
In odd characteristic a singular point would have $VW=0$.
If $W=0$, the equation gives $U=0$ and the $W$-partial is nonzero.
If $V=0$, the equation gives $U=0$ or $W^2=U^2$; the $U$-partial
is respectively $W^2$ or $-2U^2$, again nonzero. Thus the cubics
are smooth in every case, including $p=3$.

Put
\[
 k_0=\overline{\mathbb F}_p,\qquad
 K=k_0(E_0),\qquad k=\overline K,\qquad E=E_0\times k.
\]
Let $P\in E(k)$ be the geometric generic point of $E_0$, and define
\[
 M=\mathcal O_E(1),\qquad c=c_1(M),\qquad
 \lambda=\mathcal O_E(P-o)\in\operatorname{Pic}^0(E).
\]
Here $\deg M=3$. The class $\lambda$ has infinite order: under the
elliptic identification $E\simeq\operatorname{Pic}^0(E)$, finite
order would put $P$ in the finite group scheme $E_0[n]$ for some
$n>0$. Every geometric point of that finite scheme is defined over
$k_0$, contrary to genericity. Finiteness holds for all $n$, including
multiples of $p$, by
\citet[Proposition 8.1 and Theorem 8.2, pp.~12--14]{D_MilneAbelian2022}.
The actual line bundle
$\mathcal O_{E\times E}(\Delta-\{o\}\times E)$, with the second
factor as parameter, has fibre $\mathcal O_E(z-o)$ at $z$ and trivial
fibre at $o$. It will supply every ordinary family used below.

\emph{Step 2: explicit disjoint sections and a transverse finite base.}
Choose distinct constants $b_0,\ldots,b_{r-1}\in k_0$ satisfying
\[
 \begin{cases}
 b_i^2+b_i\ne0,&p=2,\\
 4-27b_i^4\ne0,&p\ne2.
 \end{cases}
\]
This is a fully specified finite choice: enumerate finite fields
$\mathbb F_{p^e}$ by increasing $e$, represent each with the first
monic irreducible polynomial in coefficient-lexicographic order,
and list its elements in the corresponding coefficient order.
Use the first field containing at least $r$ eligible elements, and
take its first $r$ eligible elements. Fix an embedding of this finite
field in $k_0$. The search terminates because only finitely many
elements are excluded. This choice affects the concrete equation but
none of its asserted properties.

Let $\ell_i=V-b_iW\in H^0(E,M)$ and let $D_i$ be its zero divisor.
The line has no intersection with $E$ at infinity: at $W=0$ it would
also have $V=0$, forcing $U=0$. In the chart $W=1$, the intersection
polynomial is $U^3+b_i^2+b_i$ in characteristic two and
$U^3-U-b_i^2$ otherwise. The first is separable precisely when
$b_i^2+b_i\ne0$; the second has discriminant $4-27b_i^4$.
Thus each $D_i$ consists of three distinct geometric points. Distinct
horizontal lines intersect only at $(1:0:0)$, which is not on $E$,
so the $D_i$ are pairwise disjoint.

Set
\[
 W_0=\mathbb P_y^{r-1}\times E,\qquad
 A=\mathcal O_{\mathbb P_y^{r-1}}(1)\boxtimes M,
 \qquad a=c_1(A)=h_y+c,
\]
and take the $r$ sections $s_i=y_i\ell_i$ of $A$.
Their common zero scheme is exactly
\[
 B=\coprod_{i=0}^{r-1}\{e_i\}\times D_i,
\]
where $e_i$ is the $i$th coordinate vertex of $\mathbb P_y^{r-1}$.
Indeed at a point of $E$ at most one $\ell_i$ vanishes; all other
coordinates $y_j$ are then forced to zero. At a point over $D_i$,
use $y_i=1$. The $r-1$ functions $y_j\ell_j$, $j\ne i$, have
independent nonzero linear terms in the $y_j$ directions, and
$\ell_i$ is a parameter in the elliptic direction. Hence the
$ds_i$ are independent and $B$ is reduced of length $3r$.
In particular no transversality assertion is inferred from a generic
choice or from a finite sample.

\emph{Step 3: the smooth incidence divisor and its regular zero scheme.}
Define
\[
 V_0=\mathbb P_x^{r-1}\times W_0,\qquad
 L=\mathcal O_{\mathbb P_x^{r-1}}(1)\boxtimes A,
 \qquad h=h_x,
\]
and let $j:Q\hookrightarrow V_0$ be the section of $L$ with equation
\[
                    f=\sum_{i=0}^{r-1}x_i y_i\ell_i=0.
\]
Off $B$, projection to $W_0$ is a $\mathbb P^{r-2}$-bundle.
Along $B$, the nonzero linear combination $\sum x_i ds_i$ cannot
vanish because the $ds_i$ are independent. Thus $Q$ is smooth.
It contains $D=\mathbb P_x^{r-1}\times B$.

Let $\mathcal U$ denote the rank-$(r-1)$ universal quotient bundle
on $\mathbb P_x^{r-1}$, with exact sequence
\[
 0\longrightarrow\mathcal O(-1)\longrightarrow
       \mathcal O^{\oplus r}\longrightarrow\mathcal U
                                                    \longrightarrow0.
\]
The map $\mathcal O^{\oplus r}\to A$ given by the $s_i$ kills
the tautological line on $Q$, precisely by the equation $f=0$.
It therefore defines a section of
$\mathcal U^\vee\otimes A$ on $Q$. Its zero scheme is $D$,
scheme-theoretically. On a chart $x_i=1$, the equation of $Q$
eliminates $s_i$ from the $r$ regular parameters of $B$, and the
remaining $r-1$ parameters generate the zero ideal. Thus the section
is regular and $\operatorname{codim}_Q D=r-1$.

The integral top-Chern formula for a regular section gives
\[
 [D]=c_{r-1}(\mathcal U^\vee\otimes A)
     =\sum_{j=0}^{r-1}(-h)^j a^{r-1-j}
                                      \quad\text{in }\mathrm{CH}^{r-1}(Q).
\]
Here $c_j(\mathcal U)=h^j$ follows from its exact sequence, and
the twist formula has coefficient one in top degree. The exact
primary statements are
\citet[Lemma 42.44.1 and Lemma 42.39.1]{D_StacksCriticalChern}.
The pullback of $\lambda$ is trivial on every component of $D$,
because its map to $E$ is constant. Projection formula consequently
gives the integral identity
\begin{equation}\label{eq:C3ay-chern-annihilator}
       \left(\sum_{j=0}^{r-1}(-h)^j a^{r-1-j}\right)\lambda=0
                                                 \quad\text{in }\mathrm{CH}^r(Q).
\end{equation}
Products with $\lambda$ here and below mean products with its first
Chern class; pullbacks of all displayed classes are understood.

\emph{Step 4: the exact source quotient and its infinite-order class.}
The integral projective-bundle formula gives
\[
 \mathrm{CH}^*(V_0)=
       \mathrm{CH}^*(E)[h,h_y]/(h^r,h_y^r).
\]
Since $c_1(L)=h+h_y+c$, eliminate $h=-(h_y+c)$ in the quotient
by this class. The relation $h^r=0$ becomes
$h_y^r+r h_y^{r-1}c=0$, because $c^2=0$ on the curve; the other
relation is $h_y^r=0$. Its degree-$r$ part is therefore exactly
\begin{equation}\label{eq:C3ay-source-quotient}
 \frac{\mathrm{CH}^r(V_0)}{c_1(L)\mathrm{CH}^{r-1}(V_0)}
       \simeq\frac{\operatorname{Pic}(E)}{\mathbb Z\cdot M^{\otimes r}}.
\end{equation}
The identification sends $h_y^{r-1}\xi$ to $\xi$.
There are no additional degree-$r$ relations: the quotient ring is
obtained by this linear elimination, and both remaining homogeneous
relations have degree $r$. The projective-bundle statement itself is
\citet[Lemma 42.36.2, Tag 02TX]{D_StacksCriticalChern}.

The subgroup $\operatorname{Pic}^0(E)$ injects into the right-hand
quotient, since a degree-zero power $M^{\otimes rn}$ has
$3rn=0$ in $\mathbb Z$, hence $n=0$. Also
$a^{r-1}\lambda=h_y^{r-1}\lambda$ because every term containing
$c\lambda$ is zero on the curve. Consequently
\[
               \gamma_0=r a^{r-1}\lambda\in\mathrm{CH}^r(V_0)
\]
has infinite order in \eqref{eq:C3ay-source-quotient}.
It is ordinarily algebraically trivial: pull the actual diagonal
family in Step~1 to $V_0\times E$, apply the fixed Chern operator
$r a^{r-1}$, and compare its fibres at $P$ and $o$.
No scalar coefficient in the field replaces this integer multiplication.

\emph{Step 5: actual strict-local realization and integral vanishing.}
Let
\[
 S=\bigoplus_{n\ge0}H^0(V_0,L^n),\qquad
 R=S_{S_+}^{\mathrm{sh}}
\]
at the geometric vertex with residue field $k$, and use the degree-one
element $f$ of Step~3 as the element of $R$.
The plane cubic has its normal degree-one full section ring: its
homogeneous ring is a normal hypersurface ring and the projective
cohomology sequence identifies it with the full section ring.
Tensoring section multiplications with the two projective-space
factors gives degree-one generation for $S$. The full section ring
is normal, equivalently by its divisorial valuation-intersection
description. Thus C.3e applies to this precise section ring.

The punctured affine cone is the punctured line bundle $L^{-1}$ over
$V_0$, hence is smooth. Its local ring at the vertex is excellent
normal of dimension $\dim V_0+1=2r$, and strict henselization retains
these properties. Its puncture $X$ is regular. The element $f$ is
nonzero and hence a nonzero divisor; its zero scheme off the vertex
is the punctured line bundle over $Q$, so $Y$ is regular too.
The map from the full punctured cone to $V_0$ is flat, and so are
the localization and strict-henselization maps; the analogous maps
over $Q$ are their base changes.

Let $\gamma$ be the pullback of $\gamma_0$ to $X$.
It is an actual ordinary class by pulling its parameter family before
taking the two fibres. The integral strict-cone injection C.3e and
\eqref{eq:C3ay-source-quotient} prove that $\gamma$ has infinite
order. On $Y$, the pullback of $L$ is trivial: the punctured line
bundle has its tautological nowhere-zero section, and this remains
true after the flat strict-local base change. Therefore $h+a=0$
on $Y$. Pulling back \eqref{eq:C3ay-chern-annihilator} gives
\[
 0=\left(\sum_{j=0}^{r-1}a^{r-1}\right)\lambda
                         =r a^{r-1}\lambda=\gamma|_Y.
\]
This is an equality in the integral Chow group in every characteristic.
In particular one does not divide by $r$ when $p\mid r$; the
non-torsion source class was explicitly chosen to be $r$ times the
elliptic class. This proves the theorem.
\end{proof}

Together with C.3aw, this proves failure of rational ordinary
injectivity in every positive characteristic throughout
$r+1\le N\le2r$ for $r\ge2$. The higher interval $N\ge2r+1$
is not settled by this construction. As in C.3x, the infinite-order
point is over the specified transcendental algebraically closed field;
ordinary rational classes over $\overline{\mathbb F}_p$ itself
vanish and are not counterexamples.
